% Proposition and proof for inclusion in the paper's theory section.
% No preamble; paste into a document that loads amsthm, amsmath, amssymb.

\begin{definition}[History Binning Key]
Let $\mathcal{O} \subseteq [0,1]^d$ be a compact observation space and let $B \in \mathbb{N}$.
The \emph{$B$-bin quantile key} $\kappa_B(h)$ of a history
$h = (o_1, a_1, \ldots, o_t, a_t)$ is obtained by replacing each observation
$o_s[0]$ (its first coordinate) with its equal-frequency quantile bin index
$\mathrm{bin}_B(o_s[0]) \in \{1, \ldots, B\}$ determined from the empirical marginal
distribution of the first observation coordinate over $N$ trajectories.
Histories that share the same key are pooled into a common \emph{history class};
their $k$-step continuation vectors are accumulated as raw $\mathbb{R}^{d_{\mathrm{eff}}}$
observations, where $d_{\mathrm{eff}} = d_{\mathcal{O}} k + d_{\mathcal{A}}(k-1)$.
\end{definition}

\begin{proposition}[History Binning Approximation Error]
\label{prop:binning-approx}
Let $\mathcal{O} \subseteq [0,1]^d$ be a compact observation space.
Let $B$ be the number of equal-frequency quantile bins applied to the first
observation dimension for history key generation.
Let $Q^*$ be the true Myhill--Nerode quotient with $n^*$ equivalence classes, and let
$\widehat{Q}_B$ be the quotient recovered by the $B$-bin algorithm on $N$ trajectories.

For each class $i \in \{1,\ldots,n^*\}$, let
$\mu_i = \mathbb{E}[o_t[0] \mid h \in \text{class } i]$
denote the mean of the first observation coordinate within that class.
Suppose the following conditions hold:
\begin{enumerate}[\upshape(a)]
    \item \textbf{Minimum signal gap.}
        $\Delta := \min_{i \neq j}\,|\mu_i - \mu_j| > 0$.
    \item \textbf{Bounded noise.}
        The marginal distribution of $o_t[0]$ within each class is sub-Gaussian with
        parameter $\sigma$, satisfying $3\sigma < \Delta/2$
        (equivalently, classes are at least $6\sigma$ apart).
    \item \textbf{Sufficient bins.}
        $B \geq 2n^*$.
\end{enumerate}
Then
\[
    \Pr\!\bigl[\widehat{Q}_B = Q^*\bigr] \;\longrightarrow\; 1
    \quad\text{as } N \to \infty.
\]
\end{proposition}

\begin{proof}[Proof sketch]
The argument proceeds in two steps.

\medskip
\noindent\textbf{Step~1: Each true class occupies at least two bins.}
Because the equal-frequency quantile boundaries track the empirical marginal distribution
of $o_t[0]$, the Glivenko--Cantelli theorem implies that these boundaries converge
uniformly to the population quantile function as $N \to \infty$.
In the population each of the $n^*$ classes occupies a positive-measure region of that
marginal (they have distinct means and sub-Gaussian tails), so for large $N$ the quantile
grid assigns each class at least $\lfloor B / n^* \rfloor \geq 2$ bins with probability
approaching~$1$.

\medskip
\noindent\textbf{Step~2: Boundary bins become single-class as $N \to \infty$.}
A \emph{boundary bin} straddles the population boundary $c_{ij} = (\mu_i + \mu_j)/2$
between adjacent classes $i$ and $j$.
An observation from class~$i$ falls in class~$j$'s bin only if it deviates from $\mu_i$
by at least $\Delta/2 - w/2$, where $w$ is the bin width.
Since $\Delta/2 > 3\sigma$ (condition~(b)), a sub-Gaussian tail bound gives
\[
    \Pr[\text{obs from class } i \text{ lands in class } j\text{'s region}]
    \;\leq\; 2\exp\!\Bigl(-\frac{\Delta^2}{8\sigma^2}\Bigr)
    \;\leq\; 2e^{-9/2} < 0.022.
\]
The expected number of cross-contaminating observations in any boundary bin is
$O(N \cdot e^{-\Delta^2/(8\sigma^2)})$, while the total count in the bin grows as
$\Theta(N/B)$.
By the Borel--Cantelli lemma, the fraction of boundary-bin observations from the
wrong class converges to~$0$ almost surely.

\medskip
\noindent\textbf{Combining the steps.}
For large $N$, each true class $i$ is represented by at least two bins whose
continuation pools are exclusively populated by observations from class~$i$.
The $B$-bin key $\kappa_B(h)$ therefore acts as a refinement of the true class label:
histories in the same true class receive the same bin sequence; histories in different
true classes receive different bin sequences (their means are $\Delta$-separated).
The downstream pooling and equivalence test correctly identifies these groups, yielding
$\widehat{Q}_B = Q^*$.
\end{proof}

\begin{corollary}[K-Cycle Environment]
\label{cor:kcycle}
For the $K$-cycle environment with observation noise $\sigma \leq \Delta/(6)$
where $\Delta = 1/K$ is the uniform gap between class means, setting $B = K$
satisfies conditions~\emph{(a)}--\emph{(c)} of Proposition~\ref{prop:binning-approx}
asymptotically.
Consequently, $\Pr[\widehat{Q}_K = Q^*] \to 1$ as $N \to \infty$.
Setting $B = 2K$ provides the full two-bin slack and is recommended for finite~$N$.
\end{corollary}

\begin{remark}[Failure modes]
When $B < 2n^*$, some pairs of true classes must share bins, producing mixed
continuation pools that cause phantom classes or spurious merges in the equivalence test.
When $3\sigma \geq \Delta/2$, boundary bins retain a constant fraction of
cross-contaminating observations at any $N$, and exact recovery fails.
In both cases the correct remedy is to use the MMD-based estimator
(Section~\ref{sec:mmd-estimator}), which operates on full continuation vectors
and is not susceptible to one-dimensional quantile-boundary artifacts.
\end{remark}
